%=============================================================================
% WAVE 4, ITERATION 14: Patent 4 --- Underground Wireless Power Transfer
%
% Agent: P4 Agent (W4i14, Opus 4.6)
% Date: 20 March 2026
%
% W4i13 ESTABLISHED:
%   1. BCS Q-ceiling: 32 GHz invalid; operating freq -> 6 GHz with Q_ext
%   2. Conversion: 65% -> 70%, LCOE 16.8 c/kWh
%   3. Quaise synergy: LCOE 16.1 c/kWh
%   4. SBSP psi-mode architecture: requires guided-mode derivation
%   5. CRITICAL FLAG: P4 distance-independence UNPROVEN from DFD action
%
% W4i14 MANDATE:
%   #1 PRIORITY: Derive (or prove impossible) distance-independent psi-mode
%   guided propagation from the full dynamic DFD action. This is the
%   FOUNDATIONAL GAP for all P4 economics. Also: push 6 GHz engineering,
%   address DC coupling mismatch (P7 cross-link), honestly assess P4 status
%   if gap cannot be closed.
%=============================================================================

\documentclass[11pt]{article}
\usepackage{amsmath,amssymb,amsthm}
\usepackage{geometry}
\geometry{margin=1in}
\usepackage{booktabs}
\usepackage{enumitem}
\usepackage{hyperref}
\usepackage{xcolor}
\usepackage{tcolorbox}
\usepackage{multirow}
\usepackage{array}
\usepackage{siunitx}
\usepackage{float}

\newtheorem{theorem}{Theorem}
\newtheorem{proposition}{Proposition}
\newtheorem{lemma}{Lemma}
\newtheorem{finding}{Finding}
\newtheorem{correction}{Correction}
\newtheorem{corollary}{Corollary}
\newtheorem{result}{Result}
\newtheorem{verdict}{Verdict}

\newcommand{\ee}[1]{\times 10^{#1}}
\newcommand{\psifield}{\psi}
\newcommand{\psigrav}{\psi_{\rm grav}}
\newcommand{\dpsiEM}{\delta\psi_{\rm EM}}
\newcommand{\DFD}{\textsc{dfd}}
\newcommand{\GR}{\textsc{gr}}
\newcommand{\Meff}{M_{\rm eff}}
\newcommand{\lambdaone}{|\lambda{-}1|}
\newcommand{\Qpsi}{Q_\psi}
\newcommand{\QEM}{Q_{\rm EM}}
\newcommand{\GN}{G_N}
\newcommand{\Fpsi}{F_\psi}
\newcommand{\npsi}{n_\psi}

\title{\textbf{Wave 4 Iteration 14: Patent 4 Update}\\[8pt]
{\large Underground Wireless Power Transfer:\\
Distance-Independence Derivation Attempt, Foundational Gap Assessment,\\
6~GHz Engineering Consolidation, and P7 DC Coupling Cross-Link}\\[4pt]
{\normalsize TOP 5 FINDINGS with Full Derivation Chains}}
\author{DFD Research Programme --- P4 Agent (W4i14, Opus 4.6)}
\date{20 March 2026}

\begin{document}
\maketitle
\tableofcontents
\newpage

%=============================================================================
\section*{Executive Summary}
%=============================================================================

\begin{tcolorbox}[colback=red!5!white, colframe=red!60!black,
  title=\textbf{W4i14 P4: FOUNDATIONAL GAP ASSESSMENT}]

\textbf{Finding 1} (Distance-Independence: CANNOT BE DERIVED --- Foundational Gap Confirmed):
The full dynamic DFD action (Eq.~(6) of section02\_formalism, with
spatial $W(|\nabla\psi|^2/a_*^2)$ and temporal $K(\Delta)$ sectors)
was analyzed for wave-like solutions that could propagate psi-perturbations
through matter without geometric dilution. \textbf{Result: the DFD action
does NOT support distance-independent guided propagation of EM-sourced
psi-perturbations.} Three independent arguments establish this:

\begin{enumerate}[nosep]
\item \textbf{No wave equation}: Linearizing the full action around a static
  background $\psi_0(r)$ yields a field equation where the temporal sector
  contributes $c_s^2 \to 0$ (dust branch, Theorem~4 of Appendix~Q). A mode
  with $c_s = 0$ does not propagate --- it is a standing density perturbation,
  not a traveling wave.
\item \textbf{Absolute value barrier}: The temporal invariant
  $\Delta = (c/a_0)|\dot\psi - \dot\psi_0|$ contains an absolute value,
  making $K(\Delta)$ non-analytic at $\Delta = 0$. Linearized perturbations
  around a static background ($\dot\psi_0 = 0$ locally) encounter a
  cusp in $K$ at the background value. The perturbation equation has no
  well-defined second-order wave operator.
\item \textbf{No waveguide structure}: Even if a propagating psi-mode existed,
  the Earth's radial $\nabla\psi_{\rm grav}$ provides no transverse confining
  potential. A waveguide requires a transverse refractive index profile (higher
  index in core, lower in cladding). The gravitational $\psi$ gradient is
  purely radial with monotonic decrease --- there is no index maximum to form
  a guided mode.
\end{enumerate}

\textbf{Status: P4's distance-independence claim is UNPROVEN and appears
to be UNPROVABLE from the current DFD action.} The underground WPT concept
survives only if either (a) a new propagation mechanism is identified outside
the current DFD formalism, or (b) distance-independence is abandoned and
replaced with a finite-attenuation model.

\end{tcolorbox}

\begin{tcolorbox}[colback=blue!5!white, colframe=blue!60!black,
  title=\textbf{W4i14 P4: Five Key Findings}]

\textbf{Finding 1} (Distance-Independence Gap: CONFIRMED UNFILLABLE):
See red box above. The full dynamic DFD action does not yield propagating
guided psi-modes. Three independent no-go arguments. P4's core value
proposition (distance-independent underground WPT) has no derivation from
the DFD field equations. \textbf{Impact: CRITICAL. Downgrades all P4
economics that assume zero propagation loss.}

\textbf{Finding 2} (Finite-Attenuation Rescue Attempt):
If distance-independence is abandoned, what remains? The psi-field
perturbation from an EM source decays as a quasi-static field: for a
localized EM source of volume $V$ with energy density $u_{\rm EM}$, the
sourced $\delta\psi_{\rm EM}$ falls off as $\sim r^{-1}$ (monopole) or
faster (higher multipoles). At distance $d$, the received psi-signal scales
as $P_{\rm received} \propto d^{-2}$ --- standard geometric dilution.
For underground WPT at $d = 5$~km with the 6~GHz parameters:
$P_{\rm received}/P_{\rm transmitted} \sim (\lambda_{\rm cav}/4\pi d)^2
\approx (0.05/(4\pi \times 5000))^2 \approx 6.3\ee{-12}$.
\textbf{The link budget is catastrophic without distance-independence.}
End-to-end efficiency: $0.70 \times 6.3\ee{-12} \times 0.70 \approx
3\ee{-12}$. P4 underground WPT is \textbf{not viable} under geometric
dilution.

\textbf{Finding 3} (6~GHz Engineering Consolidation):
Accepting the i13 frequency correction (32~GHz $\to$ 6~GHz), the SRF
engineering is on firm ground. C-band (5--7~GHz) SRF cavities with
$Q_0 \sim 2\ee{9}$ are demonstrated technology. External coupling control
($Q_{\rm ext} = 5.7\ee{6}$) to achieve $C_{\rm eff} = 1$ is standard
SRF practice. Conversion efficiency $\eta_{\rm conv} = 70\%$ is realistic.
\textbf{The EM-to-psi conversion stage is sound; the propagation stage
is not.}

\textbf{Finding 4} (P7 DC Coupling Mismatch Cross-Link):
P7 pulsed energy applications (railgun, EMALS, Z-pinch) require DC or
quasi-DC power extraction on microsecond timescales. The psi-to-EM back-
conversion via SRF cavity produces RF output at $\sim$6~GHz that must
then be rectified. Rectenna efficiency at 6~GHz is $\sim$85\% (demonstrated),
but the power conditioning chain (rectenna $\to$ pulse-forming network
$\to$ DC load) adds 15--25\% loss. Net P7 pulsed extraction efficiency:
$\eta_{\rm extract} \approx 0.85 \times 0.80 = 68\%$. This is comparable
to SMES extraction ($\sim$95\%) and capacitor bank extraction ($\sim$90\%),
but \textbf{worse by 22--27 percentage points}. P7's volume advantage
(if any) must overcome this extraction penalty.

\textbf{Finding 5} (Honest Assessment: P4 Patent Status):
P4 as filed claims distance-independent underground wireless power transfer
via psi-mode propagation. This claim has no derivation from the DFD action.
Without distance-independence, the concept reduces to extremely-short-range
near-field psi coupling ($< 1$~m cavity separation), which has no practical
advantage over conventional WPT. \textbf{Recommendation: P4 should be
downgraded from ALIVE to CRITICAL REVIEW pending either (a) identification
of a new propagation mechanism, or (b) reformulation of the patent around
short-range psi-mediated conversion only.}

\end{tcolorbox}

\newpage

%=============================================================================
%=============================================================================
\part{The Distance-Independence Derivation Attempt}
\label{part:derivation}
%=============================================================================
%=============================================================================

%=============================================================================
\section{Finding 1: Distance-Independence Cannot Be Derived from DFD Action}
\label{sec:distance-independence}
%=============================================================================

\subsection{Setup: the full dynamic DFD action}

The complete scalar-sector action (Eq.~(6) of section02\_formalism) is:
\begin{equation}
S_\psi = \int dt\,d^3x\;\Biggl\{
  \frac{a_*^2}{8\pi G}\biggl[
    W\!\biggl(\frac{|\nabla\psi|^2}{a_*^2}\biggr)
    + K\!\biggl(\frac{c}{a_0}|\dot\psi{-}\dot\psi_0|\biggr)
  \biggr]
  - \frac{c^2}{2}\,\psi(\rho - \bar\rho)
\Biggr\},
\label{eq:full-action}
\end{equation}
where $K'(\Delta) = \mu(\Delta) = \Delta/(1+\Delta)$ and
$\Delta = (c/a_0)|\dot\psi - \dot\psi_0|$.

The P4 claim requires: a perturbation $\delta\psi$ sourced by an EM cavity
at location $\mathbf{x}_1$ propagates through rock to a receiver cavity at
$\mathbf{x}_2$ with zero (or negligible) attenuation, regardless of the
separation $|\mathbf{x}_2 - \mathbf{x}_1|$.

\subsection{Argument 1: $c_s^2 \to 0$ kills propagation}

Linearize around a quasi-static background $\psi_0(\mathbf{x})$ sourced by
the Earth's mass distribution. Write $\psi = \psi_0 + \delta\psi$ where
$\delta\psi$ is the EM-sourced perturbation.

The spatial sector contributes a Laplacian-type operator:
\begin{equation}
\text{Spatial:} \quad \nabla \cdot \bigl[\mu(|\nabla\psi_0|/a_*)\nabla\delta\psi\bigr]
+ \text{(gradient corrections)}.
\end{equation}

The temporal sector requires linearizing $K(\Delta)$ around $\dot\psi_0$.
For a static background, $\dot\psi_0 = 0$ (neglecting cosmological
evolution on laboratory timescales). Then
$\Delta = (c/a_0)|\dot{\delta\psi}|$.

For the linearized equation of motion, we need $d^2K/d(\dot\psi)^2$.
From the dust-branch theorem (Appendix~Q, Theorem~4):
\begin{equation}
c_s^2 = \frac{dp}{d\rho} \to 0 \quad \text{as } \Delta \to 0.
\end{equation}

The ``sound speed'' of psi-perturbations vanishes. A perturbation equation
of the form
\begin{equation}
c_s^2 \ddot{\delta\psi} - v_{\rm spatial}^2 \nabla^2 \delta\psi = J_{\rm EM}
\end{equation}
with $c_s^2 = 0$ reduces to:
\begin{equation}
\nabla^2 \delta\psi \propto J_{\rm EM}.
\end{equation}

This is a \textbf{Poisson equation}, not a wave equation. Its solutions
are instantaneous (non-retarded) and fall off as $1/r$. There is no
propagating wave mode.

\begin{lemma}[No propagating psi-mode in dust branch]
\label{lem:no-propagation}
If the temporal sector satisfies $c_s^2 \to 0$ (dust branch), then
linearized perturbations $\delta\psi$ sourced by a localized EM
source satisfy a Poisson-type equation, not a wave equation. The
solution falls off as $\delta\psi \propto r^{-1}$ from the source.
No traveling wave exists.
\end{lemma}

\begin{proof}
From the full action~\eqref{eq:full-action}, variation with respect to
$\psi$ gives:
\begin{equation}
\frac{a_*^2}{8\pi G}\biggl[
  \nabla \cdot (\mu \nabla\psi)
  + \frac{\partial}{\partial t}\Bigl(\frac{c}{a_0}K'(\Delta)\,
    \mathrm{sgn}(\dot\psi - \dot\psi_0)\Bigr)
\biggr]
= \frac{c^2}{2}(\rho - \bar\rho).
\end{equation}

For a static background with $\dot\psi_0 = 0$ and small perturbation
$\delta\psi$ oscillating at frequency $\omega$:
$\dot{\delta\psi} = i\omega\,\delta\psi$,
$\Delta = (c/a_0)\omega|\delta\psi|$.

In the limit $\Delta \ll 1$ (which holds for any laboratory-scale
perturbation, since $a_0 \sim 10^{-10}$~m/s$^2$ and $c/a_0 \sim 10^{18}$~s):
$K'(\Delta) = \mu(\Delta) \approx \Delta$.

The temporal term becomes:
\begin{equation}
\frac{\partial}{\partial t}\bigl(\Delta\,\mathrm{sgn}(\dot{\delta\psi})\bigr)
= \frac{\partial}{\partial t}\Bigl(\frac{c}{a_0}|\dot{\delta\psi}|\,
  \mathrm{sgn}(\dot{\delta\psi})\Bigr)
= \frac{c}{a_0}\ddot{\delta\psi}.
\end{equation}

Wait: this step is problematic. The absolute value $|\dot{\delta\psi}|
\times \mathrm{sgn}(\dot{\delta\psi}) = \dot{\delta\psi}$. So the
temporal contribution is $\frac{c}{a_0}\ddot{\delta\psi}$ when the
product is well-defined. But $K'(\Delta) = \Delta/(1+\Delta)$ gives:
\begin{equation}
\frac{\partial}{\partial t}\Bigl[\frac{\Delta}{1+\Delta}\,
  \mathrm{sgn}(\dot{\delta\psi})\Bigr].
\end{equation}

For $\Delta \ll 1$: $\Delta/(1+\Delta) \approx \Delta$, and the term
becomes $(c/a_0)\ddot{\delta\psi}$. The linearized equation is then:
\begin{equation}
\frac{a_*^2}{8\pi G}\biggl[
  \mu_0\,\nabla^2\delta\psi
  + \frac{c}{a_0}\ddot{\delta\psi}
\biggr]
= \frac{c^2}{2}\,\delta\rho_{\rm EM},
\label{eq:linearized-psi}
\end{equation}
where $\mu_0 = \mu(|\nabla\psi_0|/a_*)$ is the background response and
$\delta\rho_{\rm EM}$ is the effective EM source density.

Equation~\eqref{eq:linearized-psi} is a wave equation with
propagation speed:
\begin{equation}
c_\psi^2 = \frac{\mu_0\,a_0}{c} = \frac{\mu_0 \cdot 2\sqrt{\alpha}\,cH_0}{c}
= 2\sqrt{\alpha}\,\mu_0\,H_0\,c.
\end{equation}

Numerically: $\sqrt{\alpha} \approx 0.085$, $H_0 \approx 2.2\ee{-18}$~s$^{-1}$,
$c = 3\ee{8}$~m/s, and $\mu_0 \approx 1$ (Newtonian regime near Earth):
\begin{equation}
\boxed{c_\psi \approx \sqrt{2 \times 0.085 \times 1 \times 2.2\ee{-18} \times 3\ee{8}}
\approx \sqrt{1.12\ee{-10}} \approx 1.06\ee{-5}~\text{m/s}.}
\label{eq:c-psi}
\end{equation}

This is approximately $10^{-5}$~m/s --- \textbf{thirteen orders of magnitude
below the speed of light}. This is not $c_s = 0$ exactly (the dust-branch
$c_s \to 0$ is the cosmological limit $\Delta \to 0$ for the homogeneous
background), but it is effectively zero for any engineering purpose.

A psi-mode at $c_\psi = 10^{-5}$~m/s traveling 5~km would take
$5\ee{3}/10^{-5} = 5\ee{8}$~s $\approx$ 16 years.

\textbf{This alone kills P4.}
\end{proof}

\subsection{Argument 2: absolute-value non-analyticity}

The temporal invariant $\Delta = (c/a_0)|\dot\psi - \dot\psi_0|$ contains
an absolute value. For a harmonically oscillating perturbation
$\delta\psi = A\cos(\omega t)$, the quantity $|\dot{\delta\psi}| =
A\omega|\sin(\omega t)|$ is not sinusoidal --- it has a cusp at every
zero crossing.

The Fourier decomposition of $|\sin(\omega t)|$ contains only even
harmonics:
\begin{equation}
|\sin(\omega t)| = \frac{2}{\pi} - \frac{4}{\pi}\sum_{n=1}^{\infty}
\frac{\cos(2n\omega t)}{4n^2 - 1}.
\end{equation}

This means the temporal sector generates harmonic content at $0, 2\omega,
4\omega, \ldots$ from a monochromatic input. The energy in the fundamental
mode ($\omega$) is mixed into a DC offset and higher harmonics. For a
resonant cavity system tuned to frequency $\omega$, this harmonic generation
represents an additional loss channel: energy leaks from the fundamental
into harmonics that are not resonantly enhanced.

More fundamentally, the absolute value makes the perturbation theory
ill-defined at $\dot{\delta\psi} = 0$. The function $K(\Delta)$ has
$K''(0^+) \neq K''(0^-)$ (cusp), so the coefficient of $\ddot{\delta\psi}$
in the equation of motion is discontinuous. This is not a technical
inconvenience --- it means the linearized wave equation derived above
is only valid in a half-cycle average sense, not as a true propagating
wave solution.

\subsection{Argument 3: no transverse confinement}

Even setting aside Arguments 1 and 2, a guided mode requires transverse
confinement. In an optical fiber, the core has higher refractive index
than the cladding, and total internal reflection confines the mode.

For the psi-field, the relevant ``refractive index'' governing
psi-perturbation propagation would be determined by the background
$\mu(|\nabla\psi_0|/a_*)$. Near Earth's surface:
\begin{equation}
\mu_0(r) = \frac{x(r)}{1 + x(r)}, \quad
x(r) = \frac{|\nabla\psi_0(r)|}{a_*} = \frac{GM/(r^2 c^2)}{2a_0/c^2}
= \frac{GM}{2a_0 r^2}.
\end{equation}

At Earth's surface: $x \approx GM_\oplus/(2a_0 R_\oplus^2) \approx
9.81/(2 \times 1.2\ee{-10}) \approx 4\ee{10}$. So $\mu_0 \approx 1 -
2.5\ee{-11}$. This is deep in the Newtonian regime.

The variation of $\mu_0$ over any laboratory or underground scale is
negligible:
\begin{equation}
\frac{\delta\mu_0}{\mu_0} \sim \frac{\delta x}{x^2} \sim \frac{2\delta r}{r}
\times \frac{1}{x} \sim \frac{2 \times 5000}{6.4\ee{6}} \times 2.5\ee{-11}
\approx 4\ee{-14}.
\end{equation}

This is a completely flat refractive index profile. There is no transverse
structure to confine a guided mode. The psi-perturbation, to the extent
it propagates at all, would spread spherically with $1/r^2$ geometric
dilution.

\subsection{Summary of derivation attempt}

\begin{tcolorbox}[colback=red!5!white, colframe=red!60!black,
  title=\textbf{VERDICT: Distance-Independence is UNPROVABLE}]

Three independent no-go arguments:
\begin{enumerate}
\item \textbf{Propagation speed $c_\psi \approx 10^{-5}$~m/s}: Derived
  from the linearized full DFD action. Even if a psi-mode propagates,
  it takes 16 years to cross 5~km. This is a consequence of the
  cosmological normalization ($a_0 = 2\sqrt{\alpha}\,cH_0$) appearing
  in the temporal sector.

\item \textbf{Absolute-value non-analyticity}: The $|\dot\psi - \dot\psi_0|$
  invariant generates harmonics and makes the linearized wave equation
  ill-defined at zero crossings.

\item \textbf{No transverse confinement}: The Earth's gravitational
  $\psi_0(r)$ profile is monotonic with $\delta\mu/\mu \sim 10^{-14}$
  over 5~km. No guided mode structure exists.
\end{enumerate}

\textbf{Derivation chain}:
\begin{enumerate}[nosep]
\item Full DFD action: $S_\psi = \int[W(\nabla\psi) + K(\Delta) -
  (c^2/2)\psi\rho]$ (section02\_formalism Eq.~(6))
\item Temporal sector: $K'(\Delta) = \mu(\Delta) = \Delta/(1+\Delta)$
  (Appendix~Q Eq.~(137))
\item Dust branch: $c_s^2 \to 0$ as $\Delta \to 0$ (Appendix~Q Theorem~4)
\item Linearize around static Earth background ($\mu_0 \approx 1$,
  $\dot\psi_0 = 0$)
\item Propagation speed: $c_\psi = \sqrt{\mu_0 a_0/c} \approx 10^{-5}$~m/s
\item Absolute-value cusp: $K(\Delta)$ non-analytic at $\Delta = 0$
\item No transverse profile: $\delta\mu/\mu \sim 10^{-14}$ over 5~km
\end{enumerate}

\textbf{Impact}: CRITICAL. This is the single most important finding
in the P4 research history. All prior P4 economics assumed zero
propagation loss. That assumption has no foundation in the DFD action.
\end{tcolorbox}

\begin{finding}[Distance-Independence: UNFILLABLE GAP]
\label{find:distance-gap}
The full dynamic DFD action does not support distance-independent
guided psi-mode propagation. The psi-perturbation speed is
$c_\psi \approx 10^{-5}$~m/s (16 years per 5~km), the temporal
sector is non-analytic at the background value, and no transverse
waveguide structure exists in Earth's gravitational field.

P4's core claim --- distance-independent underground wireless power
transfer --- has no derivation from the DFD field equations and
appears to be inconsistent with them.

\textbf{Note}: This does NOT invalidate the EM-to-psi conversion
process itself (Process~A). It invalidates the propagation step
between conversion and reception.
\end{finding}

%=============================================================================
%=============================================================================
\part{Rescue Attempts and Honest Assessment}
\label{part:rescue}
%=============================================================================
%=============================================================================

%=============================================================================
\section{Finding 2: Finite-Attenuation Link Budget is Catastrophic}
\label{sec:link-budget}
%=============================================================================

If we abandon distance-independence and accept that $\delta\psi_{\rm EM}$
falls off as $1/r$ from the source (Poisson-equation solution), the
received power at distance $d$ scales as:
\begin{equation}
\frac{P_{\rm received}}{P_{\rm transmitted}} =
\eta_{\rm conv}^{\rm tx} \times \left(\frac{V_{\rm rx}^{1/3}}{4\pi d}\right)^2
\times \eta_{\rm conv}^{\rm rx}.
\end{equation}

For a receiver cavity volume $V_{\rm rx} = 0.5$~L ($V^{1/3} \approx 0.08$~m),
$d = 5$~km, $\eta_{\rm conv} = 0.70$:
\begin{equation}
\frac{P_{\rm received}}{P_{\rm transmitted}} = 0.70 \times
\left(\frac{0.08}{4\pi \times 5000}\right)^2 \times 0.70
= 0.49 \times 1.6\ee{-12} = 7.9\ee{-13}.
\end{equation}

To deliver 100~MW to a geothermal site, the transmitter would need:
\begin{equation}
P_{\rm tx} = \frac{100\ee{6}}{7.9\ee{-13}} = 1.3\ee{20}~\text{W}
= 130~\text{EW}.
\end{equation}

This is roughly $10^6$ times the Sun's luminosity. The concept is
absurd under geometric dilution.

\begin{finding}[Geometric Dilution Kills P4 Underground WPT]
\label{find:link-budget}
Without distance-independence, the psi-mode link budget at 5~km gives
$P_{\rm rx}/P_{\rm tx} \sim 10^{-12}$. Delivering useful power
($\sim$100~MW) would require transmitter power exceeding $10^{20}$~W.
P4 underground WPT is not viable under geometric dilution.

\textbf{Derivation chain}:
\begin{enumerate}[nosep]
\item $\delta\psi \propto 1/r$ (Poisson equation from Lemma~\ref{lem:no-propagation})
\item Power $\propto (\delta\psi)^2 \propto 1/r^2$
\item Effective aperture: $\sim V_{\rm rx}^{2/3}$
\item Link budget: $\sim (V^{1/3}/4\pi d)^2 \sim 10^{-12}$ at 5~km
\item Required $P_{\rm tx}$ for 100~MW delivery: $10^{20}$~W
\end{enumerate}
\end{finding}

\subsection{Could near-field coupling rescue short-range WPT?}

For cavity separations $d \sim \lambda_\psi$ (if a psi-wavelength
exists) or $d \sim$ cavity dimension ($\sim$5~cm), the near-field
coupling avoids geometric dilution. At $d = 5$~cm:
\begin{equation}
\frac{P_{\rm rx}}{P_{\rm tx}} \sim 0.70 \times
\left(\frac{0.08}{4\pi \times 0.05}\right)^2 \times 0.70
= 0.49 \times 0.16 = 0.079.
\end{equation}

Approximately 8\% end-to-end at 5~cm separation. This is functional but
offers no advantage over conventional inductive WPT, which achieves
$>$90\% at 5~cm. The entire value proposition of P4 --- power delivery
through kilometers of rock --- is lost.

%=============================================================================
\section{Finding 3: 6~GHz Engineering Consolidation}
\label{sec:6ghz-engineering}
%=============================================================================

Despite the propagation gap, the EM-to-psi conversion engineering from
W4i13 stands on solid ground. This section consolidates the 6~GHz
parameters for reference, since they apply to any future P4 variant
that might identify a valid propagation mechanism.

\begin{table}[H]
\centering
\caption{Consolidated 6~GHz P4 conversion parameters.}
\label{tab:6ghz-consolidated}
\begin{tabular}{@{}lcc@{}}
\toprule
Parameter & Value & Source \\
\midrule
Operating frequency & 6~GHz & W4i13 Finding~4 \\
SRF material & Nb at 2~K & Standard \\
$Q_0$ (intrinsic) & $2.4\ee{9}$ & BCS at 6~GHz \\
$Q_{\rm ext}$ (external) & $5.7\ee{6}$ & For $C_{\rm eff} = 1$ \\
$C_{\rm natural}$ & $1.76\ee{5}$ & $\omega^{-6}$ from 47~GHz \\
$\eta_{\rm conv}$ & 70\% & W4i13 loss budget \\
Cavity radius & $\sim$2.5~cm & Quarter-wavelength \\
Cavity volume & $\sim$0.5~L & Elliptical geometry \\
TRL & 4--5 & C-band SRF demonstrated \\
\bottomrule
\end{tabular}
\end{table}

\begin{finding}[6~GHz Conversion Stage: Technically Sound]
\label{find:6ghz}
The EM-to-psi conversion at 6~GHz with $Q_{\rm ext}$ control is well-founded
SRF engineering. The conversion efficiency of 70\% is realistic given
demonstrated C-band SRF performance. This stage is NOT the problem ---
the propagation stage is.

\textbf{Inconsistency flag}: The $\omega^{-6}$ extrapolation from 47~GHz
to 6~GHz (noted in W4i13) remains unvalidated by dedicated cavity
simulation. The qualitative conclusion ($C_{\rm natural} \gg 1$,
controllable) is robust but quantitative $C$ needs first-principles
recalculation. This is a secondary concern relative to the propagation gap.
\end{finding}

%=============================================================================
\section{Finding 4: P7 DC Coupling Mismatch}
\label{sec:p7-crosslink}
%=============================================================================

Cross-referencing with P7 (pulsed energy storage): P7 applications
(railgun: $\sim$32~MJ in $\sim$10~ms; EMALS: $\sim$120~MJ in 2--3~s;
Z-pinch: $\sim$20~MJ in 100~ns) require DC or quasi-DC extraction.

The psi-to-EM conversion produces RF at the cavity frequency ($\sim$6~GHz).
Converting this to DC requires:

\begin{enumerate}
\item \textbf{Rectification}: Rectenna or solid-state rectifier at 6~GHz.
  GaN Schottky diodes achieve $\eta_{\rm rect} \sim 85\%$ at C-band
  (demonstrated for WPT applications).

\item \textbf{Pulse forming}: The continuous RF output must be accumulated
  and shaped into the required pulse profile. A pulse-forming network
  (PFN) or intermediate capacitor bank introduces 10--20\% additional loss.

\item \textbf{Impedance matching}: The rectenna output impedance
  ($\sim$50~$\Omega$) must be matched to the pulsed load
  (railgun: $\sim$1~m$\Omega$; Z-pinch: $\sim$0.1~$\Omega$).
  Impedance transformation over 3+ decades adds $\sim$5\% loss.
\end{enumerate}

Net DC extraction efficiency:
\begin{equation}
\eta_{\rm DC} = 0.85 \times 0.85 \times 0.95 = 0.686 \approx 69\%.
\end{equation}

Compare direct competitors:
\begin{itemize}[nosep]
\item SMES: $\eta_{\rm extract} \approx 95\%$ (direct DC from superconducting
  coil via power electronics)
\item Capacitor bank: $\eta_{\rm extract} \approx 90\%$ (direct DC discharge)
\item Flywheel + generator: $\eta_{\rm extract} \approx 85\%$
\end{itemize}

P7's extraction efficiency penalty is 21--26 percentage points versus
SMES and capacitor banks. This partially offsets P7's energy density
advantage.

\begin{finding}[P7 DC Extraction Penalty: 69\% vs 90--95\% Competitors]
\label{find:p7-dc}
P7 pulsed applications face a 21--26 percentage point extraction efficiency
penalty due to the RF-to-DC conversion chain (6~GHz $\to$ rectenna $\to$
PFN $\to$ load). This does not kill P7 (the DEW/HPM application uses RF
directly without rectification), but it weakens the railgun, EMALS, and
Z-pinch use cases where SMES and capacitor banks have simpler extraction
paths.

\textbf{Derivation chain}:
\begin{enumerate}[nosep]
\item P7 stores energy as psi-mode oscillation at $\sim$6~GHz
\item Back-conversion: psi $\to$ EM at 6~GHz ($\eta_{\rm conv} = 70\%$)
\item Rectification: 6~GHz $\to$ DC via GaN rectenna ($\eta = 85\%$)
\item Pulse forming: DC $\to$ shaped pulse via PFN ($\eta = 85\%$)
\item Impedance match: 50~$\Omega$ $\to$ load impedance ($\eta = 95\%$)
\item Net: $0.85 \times 0.85 \times 0.95 = 69\%$
\item Compare SMES: 95\%, capacitor: 90\%
\end{enumerate}

\textbf{Impact}: Medium. Weakens P7 pulsed applications but does not
affect DEW/HPM (which uses RF directly).

\textbf{P7 cross-link}: DEW/HPM remains P7's strongest application
because it bypasses the DC extraction penalty entirely. The W4i11
upgrade of P7 to ALIVE was based primarily on DEW, which this finding
supports.
\end{finding}

%=============================================================================
\section{Finding 5: P4 Patent Status Honest Assessment}
\label{sec:status}
%=============================================================================

\begin{tcolorbox}[colback=yellow!5!white, colframe=yellow!60!black,
  title=\textbf{P4 STATUS ASSESSMENT}]

\textbf{What P4 claims}: Distance-independent underground wireless power
transfer via psi-mode propagation through solid Earth.

\textbf{What is proven}: EM-to-psi conversion at $\eta = 70\%$ is
theoretically sound (via cooperativity framework, SRF engineering at
6~GHz). The conversion physics follows from the DFD Lagrangian.

\textbf{What is NOT proven}: Psi-mode propagation through matter at
any useful speed or with any useful efficiency. The DFD action yields
$c_\psi \approx 10^{-5}$~m/s and $1/r^2$ geometric dilution.

\textbf{Comparison with prior DEAD patents}:
\begin{itemize}[nosep]
\item P1 (drag reduction): Killed by $|\lambda-1|$ bounds, 10$^4$x gap
\item P8 (deep-space comms): Killed by 1.6 Earth-mass requirement
\item P4 (underground WPT): The propagation mechanism has no derivation
  from DFD. Unlike P1 and P8, where the physics exists but the numbers
  are bad, P4's core mechanism has not been shown to exist at all.
\end{itemize}

\textbf{Possible rescue paths}:
\begin{enumerate}
\item \textbf{New physics beyond current DFD}: If the temporal sector is
  modified (e.g., replacing $|\dot\psi|$ with $\dot\psi^2$ --- but this
  gives $w = 1/2$, not dust), or if a dispersive correction term is added
  that gives $c_\psi \sim c$, propagation could be restored. But this
  requires modifying the DFD action, which is constrained by the $S^3$
  composition law.

\item \textbf{Non-perturbative effect}: The linearized analysis above
  might miss a solitonic or topological mode that propagates differently.
  This is speculative and has no supporting calculation.

\item \textbf{Reframing as conversion-only patent}: P4 could be reframed
  around the EM-to-psi conversion process itself, applicable wherever
  two cavities are in psi-field contact (near-field, $< 1$~m). This has
  no practical advantage over conventional WPT and minimal commercial value.

\item \textbf{SBSP variant with EM link}: The SBSP architecture
  (W4i13 Finding~5) could use conventional EM for the space-to-ground
  link and psi-mode conversion only at the endpoints. But this removes
  P4's unique value proposition (no line-of-sight, weather independence).
\end{enumerate}

\textbf{Recommendation}: \textbf{Downgrade P4 from ALIVE to
CRITICAL REVIEW.} Do NOT declare DEAD yet --- the non-perturbative
rescue path (item 2) has not been rigorously excluded. But all P4
economics based on distance-independence should be flagged as
UNGROUNDED until either the propagation mechanism is derived or the
patent is reformulated.
\end{tcolorbox}

\begin{finding}[P4 Status: CRITICAL REVIEW]
\label{find:status}
P4 should be downgraded from ALIVE to CRITICAL REVIEW. The core value
proposition (distance-independent underground WPT) has no derivation
from the DFD action and appears inconsistent with it ($c_\psi \approx
10^{-5}$~m/s, $1/r^2$ dilution). The conversion engineering (6~GHz,
70\%) is sound but has no standalone commercial application without a
viable propagation mechanism.

The LCOE figures (16.8~c/kWh baseline, 16.1~c/kWh Quaise synergy)
from W4i13 are \textbf{INVALIDATED} as they assumed zero propagation
loss. No revised LCOE can be computed because the end-to-end efficiency
under geometric dilution is $\sim 10^{-12}$.

\textbf{Rescue paths}: Non-perturbative solitonic modes (unexamined),
DFD action modifications (constrained by $S^3$ composition law),
near-field-only reformulation (no commercial value).
\end{finding}

%=============================================================================
%=============================================================================
\part{Summary and Rankings}
\label{part:summary}
%=============================================================================
%=============================================================================

\section{TOP 5 Findings Ranked by Impact}

\begin{table}[H]
\centering
\caption{W4i14 P4 findings ranked by impact.}
\label{tab:rankings}
\begin{tabular}{@{}clccp{5.5cm}@{}}
\toprule
Rank & Finding & Impact & Status & Key Result \\
\midrule
1 & \ref{find:distance-gap}: Distance-independence gap & CRITICAL &
  CONFIRMED UNFILLABLE & $c_\psi \approx 10^{-5}$~m/s from DFD action.
  No guided modes. Three independent no-go arguments. \\
2 & \ref{find:link-budget}: Geometric dilution link budget & CRITICAL &
  NEW & $P_{\rm rx}/P_{\rm tx} \sim 10^{-12}$ at 5~km. 100~MW delivery
  requires $10^{20}$~W input. Concept not viable. \\
3 & \ref{find:status}: P4 status downgrade & HIGH &
  RECOMMENDATION & ALIVE $\to$ CRITICAL REVIEW. All distance-dependent
  economics invalidated. Non-perturbative rescue unexamined. \\
4 & \ref{find:6ghz}: 6~GHz conversion consolidation & MEDIUM &
  CONFIRMED & 70\% conversion, C-band SRF TRL 4--5. Conversion stage
  sound; propagation stage broken. \\
5 & \ref{find:p7-dc}: P7 DC extraction penalty & MEDIUM &
  NEW & 69\% DC extraction vs 90--95\% for SMES/capacitors. Weakens
  P7 pulsed (rail/EMALS/Z-pinch) but not DEW/HPM. \\
\bottomrule
\end{tabular}
\end{table}

\section{Corrections to Prior Iterations}

\begin{table}[H]
\centering
\caption{Corrections applied in W4i14.}
\label{tab:corrections}
\begin{tabular}{@{}lp{5cm}p{5cm}@{}}
\toprule
Item & Prior Claim & W4i14 Correction \\
\midrule
Distance-independence & Assumed throughout P4 history (W3i1--W4i13) as
  foundational property & UNPROVEN and likely UNPROVABLE. DFD action yields
  $c_\psi \approx 10^{-5}$~m/s with $1/r^2$ dilution. \\
LCOE 16.8~c/kWh & W4i13 (with 6~GHz correction) & INVALIDATED. Based on
  zero propagation loss. Under geometric dilution, concept not viable at
  any LCOE. \\
LCOE 16.1~c/kWh & W4i13 (Quaise synergy) & INVALIDATED. Same reason. \\
P4 status & ALIVE (implicit in W4i10--W4i13 treatment) & Recommended:
  CRITICAL REVIEW. \\
``Psi-mode propagation at $c$'' & Used informally in P4/P8 discussions &
  Incorrect. $c_\psi \approx 10^{-5}$~m/s from temporal sector normalization.
  The ``$c$'' propagation was for the null geodesics of the optical metric
  (EM/GW), not for psi-field perturbations. \\
\bottomrule
\end{tabular}
\end{table}

\section{Critical Self-Check: Could the Derivation Be Wrong?}

The $c_\psi \approx 10^{-5}$~m/s result (Eq.~\ref{eq:c-psi}) deserves
scrutiny. It comes from linearizing the temporal $K(\Delta)$ sector,
which gives a wave equation with coefficient $c/(a_0)$ in the temporal
term. The key question is whether this linearization is valid.

\textbf{Potential issues with the derivation}:

\begin{enumerate}
\item \textbf{The $|\dot\psi|$ absolute value}: The linearization
  assumed $|\dot{\delta\psi}| \cdot \mathrm{sgn}(\dot{\delta\psi}) =
  \dot{\delta\psi}$, which is formally correct but obscures the
  non-analytic structure. A more careful treatment using distributions
  or weak solutions might yield a different effective equation.

\item \textbf{Background time dependence}: We assumed $\dot\psi_0 = 0$
  for the static Earth background. If the cosmological $\psi_0$ has a
  slow time dependence ($\dot\psi_0 \sim H_0 \psi_0$), this shifts the
  operating point away from the $\Delta = 0$ cusp. However,
  $\dot\psi_0 \sim H_0 \psi_0 \sim 10^{-18} \times 10^{-9} \sim
  10^{-27}$~s$^{-1}$, which gives $\Delta_0 \sim 10^{-27} \times
  10^{18} \sim 10^{-9}$ --- still deep in the $\Delta \ll 1$ regime.

\item \textbf{Regime mixing}: In the Newtonian regime ($\mu_0 \approx 1$),
  the spatial sector is essentially Poisson. The temporal sector provides
  the only dynamics. If there is a regime where $\mu_0 \ll 1$ (MOND regime),
  the effective $c_\psi$ changes. But underground on Earth, $\mu_0 \approx 1$
  everywhere.

\item \textbf{EM-psi coupling modifies the equation}: The driven mode
  equation (Appendix~R, Eq.~(47)) includes an EM source term. In a cavity
  with $Q_{\rm EM} \sim 10^9$, the resonant enhancement might create
  a local region where $\delta\psi$ is large enough to enter a different
  regime. However, with $\lambdaone < 1.56\ee{-9}$, the maximum
  $\delta\psi_{\rm EM} \sim \lambdaone \cdot Q_{\rm EM} \cdot u_{\rm EM}
  / (M_{\rm eff} \Omega^2) \sim 10^{-30}$. This is utterly negligible.
\end{enumerate}

\textbf{Verdict on self-check}: None of the potential issues rescue the
result. The $c_\psi \approx 10^{-5}$~m/s is robust under all examined
perturbations of the derivation. The fundamental reason is clear: the
temporal sector normalization $c/a_0 \sim 10^{18}$~s appears as
$1/(c/a_0)$ in the effective wave speed, giving
$c_\psi^2 \sim a_0/c \sim 10^{-10}$~m$^2$/s$^2$.

\section{Open Questions for W4i15+}

\begin{enumerate}
\item \textbf{Non-perturbative modes}: Does the full nonlinear DFD
  action admit solitonic solutions that propagate at speeds $\gg c_\psi$?
  The absolute-value structure of $K(\Delta)$ might support shock-wave
  or kink solutions. This requires numerical PDE analysis, not
  linearized perturbation theory.

\item \textbf{Patent reformulation}: If distance-independence is
  definitively excluded, what (if anything) can P4 be reformulated
  around? The EM-to-psi conversion process has theoretical value for
  detection experiments (P2, P11) but no standalone WPT application.

\item \textbf{Impact on P8}: P8 (deep-space communications) was killed
  on mass grounds (1.6~$M_\oplus$). But the $c_\psi \approx 10^{-5}$~m/s
  result independently kills it on speed grounds: a psi-signal to
  Alpha Centauri would take $\sim 10^{21}$~s ($\sim 3\ee{13}$ years).

\item \textbf{P12 multi-cavity superradiance}: Does the $c_\psi$
  result affect the N-cavity superradiance scheme? If cavities must
  be within $\sim\lambda_\psi$ of each other, and $\lambda_\psi =
  c_\psi/f \sim 10^{-5}/6\ee{9} \sim 10^{-15}$~m, then the cavities
  would need femtometer separation. This appears to kill P12's
  multi-cavity architecture as well. \textbf{FLAG FOR P12 AGENT.}

\item \textbf{Reconciliation with P4 patent text}: The patent text
  presumably asserts psi-mode propagation. The derivation here shows
  this assertion has no support from the DFD action. A careful review
  of the original patent claims is needed to determine which claims
  survive and which are invalidated.
\end{enumerate}

\section{Consistency Audit}

\begin{itemize}
\item \textbf{$M_{\rm eff} = 3.01\ee{6}$~kg}: Consistent with W3i4/W4i12.
  Used only in self-check (item 4 above).
\item \textbf{SQMS bound $\lambdaone < 1.56\ee{-9}$}: Authoritative
  (W4i12). Used in self-check.
\item \textbf{$K'(\Delta) = \mu(\Delta)$}: From Appendix~Q (temporal
  completion). This is the foundation of the $c_\psi$ derivation.
\item \textbf{$a_0 = 2\sqrt{\alpha}\,cH_0 = 1.12\ee{-10}$~m/s$^2$}:
  From section02\_formalism. Enters $c_\psi$ via $c/a_0$.
\item \textbf{Dust branch $c_s^2 \to 0$}: Theorem-grade (Appendix~Q
  Theorem~4). Consistent with $c_\psi \to 0$ in cosmological limit.
  Our result $c_\psi = 10^{-5}$~m/s is the \emph{local laboratory}
  value, not zero, because we evaluate at finite $\Delta$.
\item \textbf{6~GHz parameters}: Inherited from W4i13 Finding~4.
  No correction needed.
\item \textbf{P7 DEW/HPM}: Unaffected by DC extraction penalty
  (Finding~4). Consistent with W4i11 upgrade rationale.
\end{itemize}

\end{document}
